% !TeX root = ../main.tex
\section{Assessment: where we stand and what comes next}
\label{sec:assessment}

\subsection{The app against the state of the art}

The literature splits in two (Table~\ref{tab:methods}). Methods that handle real
mixes (MuScriptor, YourMT3+, source separators) output notes or stems without
fingering; methods that output tablature (TabCNN, FretNet, Noise2Fret, TART)
require a solo guitar. A modular pipeline that chains a separator, a guitar note
model and a fingering model is currently the only way to get a tab from a
commercial recording. The app is such a pipeline, built from the strongest
components with released weights, and our measurements support the choice: on
electric guitar and on separated mixes it beats the end-to-end models we could
run by 0.17--0.23 Tab F1 (Table~\ref{tab:bench-tab}).

The most direct published comparison is TART: 54.08\% end-to-end Tab F1 averaged
over clean and noisy solo-guitar sets \reported. Our closest numbers are Tab F1
0.64--0.68 on the EGDB tones and 0.65 on separated band mixes with v2 and the
tagger \measured. These are \emph{not} the same test sets, so the only safe
conclusion is that the app is in the same range as the best published system,
on harder inputs (amplified and mixed guitar).

\begin{table}[h]
\centering
\caption{Main methods by stage. ``Mix'': designed or shown to work on
multi-instrument recordings. ``Here'': whether we ran it. Reported results use
different test sets.}
\label{tab:methods}
\small
\setlength{\tabcolsep}{3pt}
\begin{tabularx}{\textwidth}{@{}>{\raggedright}p{3.0cm}>{\raggedright}p{1.7cm}>{\raggedright}p{2.6cm}c>{\raggedright}p{2.6cm}X@{}}
\toprule
Method & Stage & Output & Mix & Here & Reported (as published) \\
\midrule
HT Demucs~\cite{rouard2023htdemucs} & separation & 4 or 6 stems & yes & app fallback & 9.20\dB{} SDR, MUSDB18 \\
BS-RoFormer~\cite{lu2023bsroformer} & separation & stems & yes & app default & 9.80\dB{} SDR; SDX23 winner \\
Banquet~\cite{watcharasupat2024banquet} & separation & queried stem & yes & not yet & beats \code{htdemucs_6s} on guitar \\
GAPS model~\cite{riley2024gaps} & notes & notes & no & app default & SOTA on GuitarSet \\
Basic Pitch~\cite{bittner2022basicpitch} & notes & notes, contour & no & bundled & lightweight baseline \\
YourMT3+~\cite{chang2024yourmt3} & notes & notes + instrument & yes & benchmarked & limited on pop \\
MuScriptor~\cite{rouard2026muscriptor} & notes & notes + instrument & yes & benchmarked & 60.4 onset F1 (own test) \\
Beat This!~\cite{foscarin2024beatthis} & rhythm & beats, downbeats & yes & app default & SOTA beat F1 without DBN \\
AlignBeat~\cite{ahn2025alignbeat} & rhythm & beats, downbeats & yes & not yet & +3.6 downbeat CMLt \\
Viterbi / HMM~\cite{sayegh1989fingering,hori2016minimax} & fingering & string, fret & -- & app fallback & hand-made costs \\
MIDI-to-Tab~\cite{edwards2024miditotab} & fingering & string & -- & no weights & preferred in a user study \\
Fretting-Transformer~\cite{hamberger2025fretting} & fingering & string, fret & -- & no weights & beats A* and Guitar Pro \\
Chen et al.~\cite{chen2015technique}, TENT~\cite{su2019tent} & techniques & per note & no & rules like these & TENT +14.7\% note F \\
TabCNN~\cite{wiggins2019tabcnn} & end-to-end & frets per frame & no & benchmarked & GuitarSet baseline \\
Noise2Fret~\cite{simionato2026noise2fret} & end-to-end & string, fret & no & own protocol only & beats baselines \\
TART~\cite{gupta2026tart} & hybrid & tab + techniques & no & no code & 81.35 F50 / 71.8 / 54.08 Tab F1 \\
\bottomrule
\end{tabularx}
\end{table}

\subsection{Decisions taken, and the evidence for each}

\begin{table}[h]
\centering
\caption{Design decisions in the transcription pipeline (all November 2026).}
\label{tab:decisions}
\small
\begin{tabularx}{\textwidth}{@{}>{\raggedright}p{4.3cm}X>{\raggedright\arraybackslash}p{1.5cm}@{}}
\toprule
Decision & Evidence & Section \\
\midrule
Modular pipeline, not end-to-end & Wins by 0.17--0.23 Tab F1 on electric guitar and mixes & \ref{sec:bench-tab} \\
BS-RoFormer SW as default separator, Demucs as fallback & +0.07 F50 on synthetic mixes, +0.12 on real songs; cost 34\,s against 10\,s per song & \ref{sec:separation}, \ref{sec:bench-moises} \\
Keep separator + CRNN rather than MuScriptor on the mix & Real songs: MuScriptor 0.42 on the mix against 0.71 for v2 on the stem; 1.3$\times$ real time & \ref{sec:bench-moises} \\
Beat This! final0 over small0 & Accompaniment CMLt 0.80 $\to$ 0.92; bar lines 35 $\to$ 41 of 60 & \ref{sec:rhythm} \\
v2 electric CRNN as an option, not the default & ``Metal High Gain'' 0.43 $\to$ 0.69, but EGDB has no licence & \ref{sec:notes} \\
Learned fingering tagger when installed & String accuracy +0.07 solo, +0.13 EGDB; final-tab 68 $\to$ 78\% & \ref{sec:fingering} \\
No hand-written octave correction for distortion & Made clean and crunch worse, metal no better & \ref{sec:notes} \\
Pan split for several guitars & Lead F1 0.40 $\to$ 0.85 in two-guitar mixes; single guitars unaffected & \ref{sec:separation} \\
Train in TF32, not bf16 & bf16 crashed (NaN) at step 1100 & \ref{sec:notes} \\
Separation-aware training (T1) moved down, fingering (T2) up, after BS-RoFormer & Stems cost only about 0.02 F50 against isolated guitar on synthetic mixes & \ref{sec:next} \\
\bottomrule
\end{tabularx}
\end{table}

\subsection{Failure modes, collected}

\begin{enumerate}
\item \textbf{Distorted guitar in real mixes} is the largest remaining loss:
  0.56 F50 on MoisesDB distorted guitars against 0.83 on acoustic.
\item \textbf{Two guitars in the same place} of the stereo image stay one
  track.
\item \textbf{Fingering on new players} is limited by training-data diversity.
\item \textbf{Techniques}: most bends, vibrato and slides are missed; palm
  mute, harmonics and dead notes are not detected at all.
\item \textbf{Rhythm}: compound meters, 16th triplets and 32nds; solo-guitar
  recordings without accompaniment.
\item \textbf{Same-pitch notes} cannot be represented by MT3-style
  transcribers (a reason not to adopt them as-is).
\item \textbf{No real-song benchmark}: the use case is not yet measured.
\end{enumerate}

\subsection{Open problems in the literature}

\begin{enumerate}
\item \textbf{Mix-level tablature}: no method outputs string and fret from a
  band recording. Adding string and fret tokens to a mix-trained transcriber
  such as MuScriptor would combine mix robustness with timbre-aware fingering,
  and would also fix the same-pitch limitation (a unison on two strings becomes
  two different tokens).
\item \textbf{Separation for transcription}: separators are trained for SDR,
  transcribers on clean stems; training each on the other's output is untested
  for guitar.
\item \textbf{Part separation}: splitting two guitars of the same kind needs
  spatial cues or a separator conditioned on the part.
\item \textbf{Techniques in context}: learned classifiers evaluated on
  amplified and polyphonic material.
\item \textbf{A public real benchmark}: commercial recordings with aligned
  reference tabs.
\item \textbf{Rhythm notation}: quantisation, tuplets and compound meters are
  done by rules and not evaluated anywhere.
\end{enumerate}

\subsection{What the agents are working on next}
\label{sec:next}

The detailed plan, with data and compute estimates, is in
\file{docs/transcription_roadmap.md}; the tasks are named T1--T5 there.

\begin{nextsteps}[: in order]
\begin{enumerate}
\item \textbf{Real-song benchmark} (Phase 0.3) \planned: collect 30--50 song +
  tab pairs, align, hand-check 20--30, record the current pipeline as baseline.
\item \textbf{T1, separation-aware CRNN} \wip: train on what inference sees.
  The uncommitted script \code{scripts/train/pseudo_label_moisesdb.py} builds
  training data from MoisesDB: the input is a separator's guitar stem of the
  real mix, the target the notes a teacher model (v2) transcribed from the
  song's isolated guitars (\term{pseudo-labels}), split by song into train
  (67\%), validation (8\%) and test (25\%). \code{finetune_crnn.py} gains
  \code{--extra-train} and \code{--extra-val} to mix these items in. The target
  is the 0.56 on distorted guitars in real songs.
\item \textbf{T2, fingering}: request DadaGP for pre-training (more players),
  then condition on audio (CRNN features at each onset, as TART does); retrain
  without EGDB if its authors do not allow distribution \planned.
\item \textbf{T3, per-note technique classifier} on GOAT, IDMT and Guitar-TECHS
  with tone augmentation \planned.
\item \textbf{T4, guitar tab transformer on full mixes}: fine-tune MuScriptor
  with string and fret tokens, small first, then large across the DGX Sparks;
  only after T1/T2 baselines exist \planned.
\item \textbf{T5, two guitars in the same spot}: a query-based or fine-tuned
  RoFormer separator for ``rhythm'' and ``lead''; lowest priority \planned.
\end{enumerate}
Smaller items: try Banquet and AlignBeat; palm mute and harmonics; 6/8,
16th triplets and 32nds; batch transcription of a Spotify playlist via YouTube;
a native (ONNX) separator so Android could separate; a tab editor.
\end{nextsteps}

\subsection{Conclusion}

Guitar tablature transcription moved from frame-wise CNNs on acoustic guitar to
modular systems with learned fingering and techniques, and to generative
end-to-end models; in parallel, large transcribers trained on real mixes now
handle distorted guitar but stop at notes. For transcribing a guitar part from a
song, the practical state of the art is a modular pipeline, and that is what the
app implements. A band-split separator recovers most of what the mix costs, a
note model trained on many tones handles amplified guitar on isolated
recordings, and a learned fingering model fixes most of the Viterbi's trouble
with solo lines. What remains is distorted guitar inside real mixes, the
separation of guitars that share a position, techniques beyond the easy
cases, and, before anything else, a benchmark of real songs with reference tabs
so that progress on the actual task can be measured.
